\subsection{Why not substitute tokens after generation?}
\label{sec:v1}

An earlier design first generated an unwatermarked text and then replaced selected tokens
with alternatives having the desired contrast sign. A candidate $z$ was admissible only if
its probability relative to the most likely token satisfied
$p(z)/p(z^\star)\geq e^{-\tau}$. This approach is simple and uses the same detector, but
the experiments show that it has too little capacity on the evaluated model.

Across 36 settings covering three carrier rules, three fluency budgets, and two guidance
weights, the median selected position had only one admissible token at $\tau=3$
($p(z)/p(z^\star)\geq0.05$). The best setting within a $1.35\times$ perplexity budget
reached only $0.10$ TPR at $1\%$ FPR, compared with $0.86$ for the local dgMARK
reproduction at $1.23\times$.

A sampler correction made the limitation clearer. Confidence had initially been computed
from Gumbel-perturbed logits rather than the clean model distribution. Fixing this error
reduced draft perplexity from $22.2$ to $9.5$ and mean conditional entropy from $0.655$ to
$0.319$. At the same time, the mean number of admissible alternatives fell from $2.5$ to
$1.6$ at $\tau=3$ and from $14.1$ to $5.6$ at $\tau=6$. In this setting, sharper
conditionals improve the base sampler while removing the alternatives needed for post-hoc
substitution.

This result motivates resampling rather than a universal claim that substitution must fail
for every DLM. ReTrace makes its choice while the model is already sampling and falls back
to an unconditional proposal when steering fails; it does not force a replacement after
the surrounding text has been fixed.
